\section[A. Synecdoche, New York]{\myboxi{      \BF{\TT{A}}} Synecdoche, New York}

\textit{Time limit}   2 seconds

\textit{Memory limit} 1 Gigabyte

\textit{Input} \qquad standard input

\textit{Output} \qquad standard output


\begin{quote}
	\textit{And they say there is no fate, but there is, it's what you create.}

	\textit{And even though the world goes on for eons and eons...}

	\textit{...you are only here for a fraction of a fraction of a second.}
\end{quote}

\subsection{Problem Statement}

Caden has come to understand ``fate'' in his own way. Leaving theater behind, he wants to use a string to represent his life through synecdoche. But in seeking a pure, original expression true to himself, he again falls into an infinite loop of self-reference. This time, Caden leaves the choice to you: portray his life for him.

Caden's synecdoche is a string $T$ of length $L$. We call $T$ \BF{valid} if and only if all three conditions below hold:

\begin{enumerate}
	\item $T$ consists only of lowercase English letters.
	\item For every $i = 1, 2, \ldots, K$, the character of $T$ at position $p_i$ must equal the given character $c_i$, where $1 \le p_i \le L$ and $c_i$ is a lowercase letter.
	\item For every $i = 1, 2, \ldots, N$, the given string $s_i$ must \BF{not} occur as a contiguous substring of $T$.
\end{enumerate}

Given $L$, $K$ position constraints $(p_i,c_i)$, and $N$ forbidden substrings $s_i$, count all valid strings $T$. Output the answer modulo $998\,244\,353$.

\subsection{Input}

The first line contains three integers $L$, $K$, and $N$ ($1 \le L \le 10^{18}$, $0 \le K \le 500$, $0 \le N \le 50$): the string length, the number of position constraints, and the number of forbidden substrings.

Each of the next $K$ lines contains an integer $p_i$ and a lowercase letter $c_i$ ($1 \le p_i \le L$), describing a position constraint. All $p_i$ are distinct.

Each of the next $N$ lines contains a nonempty string $s_i$, describing a forbidden substring.

The sum of the lengths of all $s_i$ does not exceed $100$; that is, $\sum_{i=1}^N \vert{}s_i\vert{} \le 100$.

\subsection{Output}

Output one integer for each test case: the number of valid strings modulo $998\,244\,353$.

\subsection{Sample 1}

\begin{minipage}{\linewidth}
	\begin{minipage}[t]{0.49\linewidth}
		\begin{lstlisting}[title=\texttt{standard input}]
3 1 1
2 b
ab
\end{lstlisting}
	\end{minipage}
	\begin{minipage}[t]{0.49\linewidth}
		\begin{lstlisting}[title=\texttt{standard output}]
650
\end{lstlisting}
	\end{minipage}
\end{minipage}

\subsection{Sample 2}

\begin{minipage}{\linewidth}
	\begin{minipage}[t]{0.49\linewidth}
		\begin{lstlisting}[title=\texttt{standard input}]
10 2 2
3 a
5 c
abc
dcabe
\end{lstlisting}
	\end{minipage}
	\begin{minipage}[t]{0.49\linewidth}
		\begin{lstlisting}[title=\texttt{standard output}]
113392371
\end{lstlisting}
	\end{minipage}
\end{minipage}


\subsection{Sample 1 Explanation}

The string has length $3$, its second character must be \TT{b}, and it must not contain \TT{ab}. Valid strings have the form \TT{?b?}; the first character can be any lowercase letter except \TT{a}. There are $25 \times 26 = 650$ such strings.

